\section{Introduction}

The use of large language models for simulating expert review raises a fundamental question: do AI-generated reviewer personas produce genuinely distinct perspectives, or do they converge on generic, model-characteristic feedback regardless of the persona specification?

This question is critical for the validity of AI-simulated review processes. If all personas produce similar feedback, the apparent diversity of a multi-reviewer panel is illusory, and the synthesis of ``multiple perspectives'' adds no value beyond a single review. Conversely, if personas can be calibrated to produce distinct, expertise-consistent feedback, multi-reviewer panels provide genuine multi-perspective assessment.

We address this question through empirical analysis of 186+ reviews generated by 45+ distinct reviewer personas across 14 papers. Our key findings:

\begin{enumerate}
    \item \textbf{Profile-based prompting produces measurable distinctness}: Pairwise rank correlations between reviewers range from $\rho = 0.05$ to $\rho = 0.80$, with no two reviewers producing identical assessments.
    \item \textbf{Three consistent blocs emerge}: Systems-focused (Zaharia, Shankar, Liang), Agent-focused (Chase, Khattab), and Human-Centered (Shneiderman, Amershi) reviewers form stable clusters across independent sessions.
    \item \textbf{Key-question prompting is the strongest calibration signal}: Providing each persona with a characteristic question (e.g., ``How was this evaluated?'' for Percy Liang) produces more distinct reviews than expertise tags or venue alignment alone.
    \item \textbf{No fine-tuning required}: All persona differentiation is achieved through prompt construction from structured profiles, making the method applicable with any capable LLM.
\end{enumerate}
